% from ega1-10-fr.tex
\oldpage[I]{196}

La détermination de $\widehat X_{\mathrm{red}}$ étant locale
(\hyperref[I.10.5.4-fr]{10.5.4}), on peut encore supposer que $X$ est un
schéma affine d'anneau noethérien~; avec les notations de
(\hyperref[I.10.8.5-fr]{10.8.5}), l'idéal $\mathfrak{T}$ des éléments
topologiquement nilpotents de $\widehat A$ est l'image réciproque par
l'application canonique
$\widehat A\to\widehat A/\widehat{\mathfrak{J}}=A/\mathfrak{J}$ du
nilradical de $A/\mathfrak{J}$
(\hyperref[0.7.1.3-fr]{0, 7.1.3}), donc $\widehat A/\mathfrak{T}$ est
isomorphe au quotient de $A/\mathfrak{J}$ par son nilradical. La première
assertion résulte donc de (\hyperref[I.10.5.4-fr]{10.5.4}) et
(\hyperref[I.5.1.1-fr]{5.1.1}). Si $\widehat X_{\mathrm{red}}$ est
noethérien, son espace sous-jacent $X'$ l'est aussi, donc les
$X'_n=\operatorname{Spec}(\mathscr{O}_X/\mathscr{J}^n)$ sont noethériens
(\hyperref[I.6.1.2-fr]{6.1.2}) et il en est de même de $\widehat X$
(\hyperref[I.10.6.4-fr]{10.6.4})~; la réciproque est immédiate, en vertu
de (\hyperref[I.6.1.2-fr]{6.1.2}).

\begin{env}[10.8.7]
\label{I.10.8.7-fr}
Les homomorphismes canoniques
$\mathscr{O}_X\to\mathscr{O}_X/\mathscr{J}$ (pour
$\mathscr{J}\in\Phi$) forment un système projectif et donnent donc, par
passage à la limite projective, un homomorphisme de faisceau d'anneaux
$\theta:\mathscr{O}_X\to\psi_*((\mathscr{O}_X)_{/X'})
=\varprojlim_\Phi(\mathscr{O}_X/\mathscr{J})$, en désignant par $\psi$
l'injection canonique $X'\to X$ des espaces sous-jacents. Nous désignerons
par $i$ (ou $i_X$) le morphisme (dit canonique)
\[
  (\psi,\theta):X_{/X'}\to X
\]
d'espaces annelés.

Par tensorisation, pour tout $\mathscr{O}_X$-Module cohérent
$\mathscr{F}$, les homomorphismes canoniques
$\mathscr{O}_X\to\mathscr{O}_X/\mathscr{J}$ donnent des homomorphismes
$\mathscr{F}\to
\mathscr{F}\otimes_{\mathscr{O}_X}(\mathscr{O}_X/\mathscr{J})$ de
$\mathscr{O}_X$-Modules qui forment encore un système projectif, et
donnent donc, par passage à la limite projective, un homomorphisme
canonique fonctoriel
$\gamma:\mathscr{F}\to\psi_*(\mathscr{F}_{/X'})$ de
$\mathscr{O}_X$-Modules.
\end{env}

\begin{proposition}[10.8.8]
\label{I.10.8.8-fr}
\medskip\noindent
\begin{enumerate}
  \item[\rm{(i)}] Le foncteur $\mathscr{F}_{/X'}$ (en $\mathscr{F}$) est
    exact.
  \item[\rm{(ii)}] L'homomorphisme fonctoriel
    $\gamma^\sharp:i^*(\mathscr{F})\to\mathscr{F}_{/X'}$ de
    $(\mathscr{O}_X)_{/X'}$-Modules est un isomorphisme.
\end{enumerate}
\end{proposition}

\medskip\noindent
\begin{enumerate}
  \item[\rm{(i)}] Il suffit de prouver que si
    $0\to\mathscr{F}'\to\mathscr{F}\to\mathscr{F}''\to0$ est une suite
    exacte de $\mathscr{O}_X$-Modules cohérents, et $U$ un ouvert affine
    de $X$, d'anneau noethérien $A$, la suite
    \[
      0\to\Gamma(U\cap X',\mathscr{F}'_{/X'})
      \to\Gamma(U\cap X',\mathscr{F}_{/X'})
      \to\Gamma(U\cap X',\mathscr{F}''_{/X'})\to0
    \]
    est exacte. On a alors $\mathscr{F}|U=\widetilde M$,
    $\mathscr{F}'|U=\widetilde{M'}$, $\mathscr{F}''|U=\widetilde{M''}$,
    où $M$, $M'$, $M''$ sont trois $A$-modules de type fini tels que la
    suite $0\to M'\to M\to M''\to0$ soit exacte
    (\hyperref[I.1.5.1-fr]{1.5.1} et
    \hyperref[I.1.3.11-fr]{1.3.11})~; soit $\mathscr{J}\in\Phi$ et soit
    $\mathfrak{J}$ un idéal de $A$ tel que
    $\mathscr{J}|U=\widetilde{\mathfrak{J}}$. On a alors
    \[
      \Gamma(U\cap X',
      \mathscr{F}\otimes_{\mathscr{O}_X}(\mathscr{O}_X/\mathscr{J}^n))
      =M\otimes_A(A/\mathfrak{J}^n)
    \]
    (\hyperref[I.1.3.12-fr]{1.3.12})~; donc, par définition de la limite
    projective, on a
    \[
      \Gamma(U\cap X',\mathscr{F}_{/X'})
      =\varprojlim_n(M\otimes_A(A/\mathfrak{J}^n))=\widehat M
    \]
    séparé complété de $M$ pour la topologie $\mathfrak{J}$-préadique, et
    de même
    \[
      \Gamma(U\cap X',\mathscr{F}'_{/X'})=\widehat{M'},\qquad
      \Gamma(U\cap X',\mathscr{F}''_{/X'})=\widehat{M''}~;
    \]
    notre assertion résulte alors de ce que lorsque $A$ est noethérien,
    le foncteur $\widehat M$ en $M$ est exact sur la catégorie des
    $A$-modules de type fini (\hyperref[0.7.3.3-fr]{0, 7.3.3}).

